\begin{table}[H]
\small
\setlength{\tabcolsep}{5pt}
\renewcommand{\arraystretch}{0.95}

\caption{\label{tab:country-heterogeneity-cumulative-gap}Cumulative inflation gap between the first and fifth income quintiles, 2020--2025}
\centering
\begin{tabular}[t]{llll}
\toprule
Country & Quintile 1 cumulative inflation & Quintile 5 cumulative inflation & Quintile 1-Quintile 5 gap\\
\midrule
Estonia & 53.9 & 46.3 & 7.6\\
Lithuania & 43.2 & 40.2 & 3.0\\
Latvia & 44.0 & 36.8 & 7.2\\
Slovakia & 37.0 & 38.1 & -1.1\\
Austria & 28.3 & 28.3 & 0.0\\
\addlinespace
Finland & 16.1 & 17.8 & -1.7\\
Luxembourg & 20.7 & 20.5 & 0.2\\
Netherlands & 27.0 & 26.8 & 0.2\\
Greece & 22.0 & 20.8 & 1.2\\
Ireland & 24.7 & 18.2 & 6.5\\
\addlinespace
Portugal & 20.9 & 20.5 & 0.4\\
Cyprus & 18.1 & 19.0 & -0.9\\
Germany & 22.6 & 25.3 & -2.7\\
France & 17.1 & 18.1 & -1.0\\
Italy & 21.7 & 21.5 & 0.2\\
\addlinespace
Spain & 21.6 & 21.7 & -0.1\\
Euro area & 22.1 & 22.8 & -0.7\\
\bottomrule
\end{tabular}
\vspace{-0.4em}
\begin{flushleft}
\footnotesize{\emph{Notes:} Entries are percentage points. Cumulative inflation is computed from average income-quintile price indices in 2020 and 2025, using RAS expenditure weights. The gap is Quintile 1 minus Quintile 5. \emph{Sources:} Eurostat; authors' calculations.}
\end{flushleft}
\end{table}
